\chapter{Matching multi-pattern: keyword tree e Aho--Corasick}
\label{cap:aho-corasick}

\scan{12}
Generalizziamo il problema del pattern matching esatto al caso in cui i pattern da cercare
non siano uno solo, ma un'intera collezione.

\section{Il problema: exact set matching}
\label{sec:exact-set-matching}

\begin{definizione}[Exact set matching]\label{def:exact-set-matching}
Dati un testo $T \in \S^{*}$ e una collezione di pattern
\[
  \mathcal{P} \;=\; \set{P_1, \dots, P_z}, \qquad P_i \in \S^{*},
\]
trovare tutte le occorrenze in $T$ di tutti gli elementi di $\mathcal{P}$.
\end{definizione}

\paragraph{Notazione.}
Poniamo $n_i = \abs{P_i}$ per $i \in \set{1,\dots,z}$ e
\[
  n \;=\; \sum_{i=1}^{z} n_i ,
\]
e indichiamo con $k$ il numero complessivo di occorrenze da riportare, cioè la dimensione
dell'output. In questo capitolo $n$ è dunque la lunghezza \emph{complessiva dei pattern} e
non quella del testo: la lunghezza del testo la scriviamo sempre, esplicitamente,
$\abs{T}$. Lo scenario tipico è quello in cui $T$ è enorme rispetto ai pattern, per cui il
costo dominante è il numero di scansioni del testo.

\paragraph{Soluzione naive.}
Applichiamo $z$ volte il miglior algoritmo per singolo pattern di cui disponiamo (KMP,
oppure l'algoritmo $Z$), che costa $\Oh(\abs{T} + n_i)$ sul pattern $P_i$. In totale
\[
  \sum_{i=1}^{z} \Oh\bigl(\abs{T} + n_i\bigr)
  \;=\; \Oh\bigl(z\,\abs{T} + n\bigr).
\]

Il \emph{punto debole} è evidente: stiamo scandendo $T$ per $z$ volte. È il fattore $z$
davanti ad $\abs{T}$ che vogliamo eliminare.

\paragraph{La sfida.}
Riusciamo a ridurci a \emph{una sola} scansione del testo? Per farlo dobbiamo combinare
tutti i pattern di $\mathcal{P}$ in un'unica struttura dati, cioè in un albero.

\section{Il keyword tree}
\label{sec:keyword-tree}

\scan{12}
Il \emph{keyword tree} $K(\mathcal{P})$ è l'albero che si ottiene sovrapponendo i prefissi
comuni degli elementi di $\mathcal{P}$: ogni cammino dalla radice a una foglia si legge come
uno dei pattern e, viceversa, ogni pattern corrisponde al cammino verso una foglia. Il
guadagno è immediato: leggendo un solo carattere di $T$ e scendendo di un arco si avanza
simultaneamente su \emph{tutti} i pattern che condividono il prefisso già letto, cioè si
scandiscono più elementi di $\mathcal{P}$ in parallelo.

Nel disegno sono già indicati, tratteggiati, i \emph{failure link}: sono i collegamenti che
permettono di riprendere la ricerca da un altro nodo dell'albero quando la discesa si blocca
su un mismatch, senza mai arretrare sul testo. Li definiremo formalmente nella
\autoref{sec:failure-link}.

\begin{figure}[H]
\centering
% Vista panoramica del keyword tree K(P): otto cammini radice-foglia (i pattern
% P_1..P_8) e due failure link tratteggiati.  I tratti in accento sono le coppie
% di sottocammini che portano la stessa stringa (suffisso/prefisso): per ogni
% failure link <v, n_v> sono marcati il cammino root ~> n_v e il suffisso di pari
% lunghezza del cammino root ~> v.  Le due coppie sono
%   alfa = root~>A  e  H2~>C   (failure link C -> A)   lunghezza 4.611
%   beta = root~>D  e  H~>B    (failure link B -> D)   lunghezza 1.392
% I due tratti di ciascuna coppia hanno la STESSA lunghezza (e' l'informazione
% del disegno).  Il tratto di destra parte da H2, sotto D, e non da D: nel
% manoscritto fra i due tratti marcati del bordo destro c'e' un segmento a penna
% sottile che li separa, senza il quale nel PDF si fondono in un unico tratto.
\begin{tikzpicture}[
    ramo/.style    = {draw=graydark, line width=0.55pt, line cap=round, line join=round},
    marcato/.style = {draw=accent, line width=2.4pt, line cap=round, line join=round,
                      opacity=0.7},
    nodo/.style    = {circle, fill=graydark, inner sep=0pt, minimum size=3.4pt},
    foglia/.style  = {circle, fill=graydark, inner sep=0pt, minimum size=2.4pt},
    flink/.style   = {draw=black, line width=0.7pt, dash pattern=on 3pt off 2.2pt,
                      -{Triangle[length=5pt,width=4pt]}},
    et/.style      = {font=\footnotesize},
  ]

  % ---- nodi -------------------------------------------------------------
  \coordinate (r)  at ( 0.00, 5.50);   % radice
  \coordinate (m1) at (-2.00, 4.07);   % biforcazione sul cammino r -> A
  \coordinate (A)  at (-3.78, 2.86);   % arrivo del failure link lungo
  \coordinate (B)  at (-0.63, 2.53);   % partenza del failure link corto
  \coordinate (D)  at ( 1.16, 4.73);   % arrivo del failure link corto (poco profondo)
  \coordinate (F)  at ( 2.73, 2.97);   % biforcazione sul bordo destro
  \coordinate (C)  at ( 4.655, 0.935); % partenza del failure link lungo (profondo)
  \coordinate (E)  at (-1.26, 1.65);
  \coordinate (G)  at ( 2.10, 1.76);
  \coordinate (P1) at (-4.83, 0.28);
  \coordinate (P2) at (-2.63, 0.28);
  \coordinate (P3) at (-1.79, 0.44);
  \coordinate (P4) at (-0.74, 0.22);
  \coordinate (P5) at ( 0.95, 0.55);
  \coordinate (P6) at ( 1.89, 0.33);
  \coordinate (P7) at ( 3.15, 0.33);
  \coordinate (P8) at ( 5.05, 0.30);
  \coordinate (H)  at ($(m1)!0.33!(B)$);    % inizio del tratto marcato che entra in B
  \coordinate (H2) at ($(D)!0.2328!(F)$);   % inizio del tratto marcato che scende in C

  % ---- tratti marcati (sotto ai rami) -----------------------------------
  \draw[marcato] (r) -- (m1) -- (A);      % alfa del failure link C -> A
  \draw[marcato] (H2) -- (F) -- (C);      % alfa del failure link C -> A
  \draw[marcato] (r) -- (D);              % beta del failure link B -> D
  \draw[marcato] (H) -- (B);              % beta del failure link B -> D

  % ---- rami dell'albero -------------------------------------------------
  \draw[ramo] (r) -- (m1) -- (A);
  \draw[ramo] (A) -- (P1);
  \draw[ramo] (A) -- (-3.10,1.95) -- (P2);
  \draw[ramo] (m1) -- (B);
  \draw[ramo] (B) -- (E);
  \draw[ramo] (E) -- (P3);
  \draw[ramo] (E) -- (-0.95,0.95) -- (P4);
  \draw[ramo] (B) -- (0.30,1.55) -- (P5);
  \draw[ramo] (r) -- (D) -- (F) -- (C) -- (P8);
  \draw[ramo] (F) -- (G);
  \draw[ramo] (G) -- (P6);
  \draw[ramo] (G) -- (2.85,1.05) -- (P7);

  % ---- pallini ----------------------------------------------------------
  \node[nodo] at (r) {};
  \node[nodo] at (A) {};
  \node[nodo] at (B) {};
  \node[nodo] at (C) {};
  \node[nodo] at (D) {};
  \foreach \l in {P1,P2,P3,P4,P5,P6,P7,P8} \node[foglia] at (\l) {};

  % ---- failure link -----------------------------------------------------
  \draw[flink] (C) .. controls (3.30,0.62) and (1.30,1.42) .. (-0.70,2.05)
                   .. controls (-1.95,2.40) and (-2.95,2.64) .. (A);
  \draw[flink] (B) .. controls (0.55,2.95) and (1.18,3.85) .. (D);
  \node[et,align=center] at (1.50,3.12) {\emph{failure}\\[-1pt]\emph{link}};

  % ---- etichette delle foglie -------------------------------------------
  \node[et,below=2pt] at (P1) {$P_1$};
  \node[et,below=2pt] at (P2) {$P_2$};
  \node[et,below=2pt] at (P3) {$P_3$};
  \node[et,below=2pt] at (P4) {$P_4$};
  \node[et,below=2pt] at (P5) {$P_5$};
  \node[et,below=2pt] at (P6) {$P_6$};
  \node[et,below=2pt] at (P7) {$P_7$};
  \node[et,below=2pt] at (P8) {$P_8$};

  \node[et,above=2pt] at (r) {$\mathrm{root}$};
\end{tikzpicture}

\caption{Il keyword tree $K(\mathcal{P})$ a colpo d'occhio: i cammini dalla radice alle $z$
foglie sono i pattern $P_1,\dots,P_z$. Tratteggiati, due failure link: puntano sempre da un
nodo a un nodo di profondità minore. In evidenza, per ciascun failure link
$\tuplev{v, \nv{v}}$, i due sottocammini che portano la stessa stringa: il cammino dalla
radice a $\nv{v}$ e il suffisso di pari lunghezza del cammino dalla radice a $v$.}
\label{fig:keyword-tree-panoramica}
\end{figure}

\subsection{Definizione}

\scan{13}

\begin{definizione}[Keyword tree]\label{def:keyword-tree}
Dato $\mathcal{P} = \set{P_1, \dots, P_z}$ con $P_i \in \S^{*}$ per
$i \in \set{1, \dots, z}$, il \emph{keyword tree} di $\mathcal{P}$ è un albero $K$ i cui
archi sono etichettati con elementi di $\S$ e tale che:
\begin{enumerate}
  \item archi uscenti dallo stesso nodo hanno etichette distinte;
  \item ogni cammino radice--foglia corrisponde a $P_i$ per qualche $i \in \set{1,\dots,z}$,
        e viceversa ogni $P_i$ si legge lungo un cammino radice--foglia.
\end{enumerate}
Scriviamo $K(\mathcal{P})$ per il keyword tree di $\mathcal{P}$.
\end{definizione}

\begin{osservazione}
Le due condizioni dicono, rispettivamente:
\begin{itemize}
  \item la (1) che ogni stringa individua in $K$ un \emph{cammino univoco}: la discesa dalla
        radice è deterministica e due nodi distinti non possono avere la stessa etichetta di
        cammino. I nodi di $K$ sono dunque in biiezione con i prefissi degli elementi di
        $\mathcal{P}$: il keyword tree è il \emph{trie} dei pattern, cioè una
        rappresentazione \emph{non ridondante sui prefissi};
  \item la (2) che i \emph{cammini} radice--foglia sono in corrispondenza con gli
        \emph{elementi di $\mathcal{P}$}.
\end{itemize}
\end{osservazione}

\subsection{Ipotesi semplificativa}

Assumiamo che nessun $P_i$ sia \emph{fattore proprio} di $P_j$ per $i \neq j$, cioè che
nessun pattern sia sottostringa propria di un altro. Scaricheremo questa ipotesi nella
\autoref{sec:output-link}.

\begin{osservazione}
L'ipotesi garantisce che foglie di $K(\mathcal{P})$ ed elementi di $\mathcal{P}$ siano in
corrispondenza biunivoca --- in particolare le foglie sono esattamente $z$ --- e soprattutto
che durante la ricerca ogni occorrenza venga scoperta arrivando in una foglia: se infatti
$P_j$ fosse sottostringa di $P_i$, percorrendo il cammino di $P_i$ si passerebbe
\emph{sopra} un'occorrenza di $P_j$ senza mai visitarne la foglia. Per la sola biiezione
foglie--pattern basterebbe in realtà l'ipotesi più debole «nessun $P_i$ è prefisso di
$P_j$».
\end{osservazione}

\subsection{Misure di \texorpdfstring{$K(\mathcal{P})$}{K(P)}}

Le tre misure di $K(\mathcal{P})$ sono
\[
\begin{aligned}
  \text{dimensione di } K(\mathcal{P}) \;&\le\; \sum_{i=1}^{z} n_i \;=\; n, \\[2pt]
  \text{altezza di } K(\mathcal{P}) \;&=\; \max_{1 \le i \le z} \set{n_i}, \\[2pt]
  \text{grado (fan-out) di } K(\mathcal{P}) \;&\le\; \min\set{\abs{\S},\, z}.
\end{aligned}
\]

\begin{osservazione}
Per \emph{dimensione} si intende il numero di archi, che in un albero è pari al numero di
nodi meno uno: gli archi di $K(\mathcal{P})$ sono tanti quanti i prefissi non vuoti
\emph{distinti} degli elementi di $\mathcal{P}$, perché ogni pattern contribuisce al più
$n_i$ archi e i prefissi condivisi vengono contati una sola volta. Dunque
$\#\text{archi} \le n$, con uguaglianza se e solo se nessuna coppia di pattern condivide un
prefisso non vuoto; lo spazio è quindi $\Oh(n)$, e $\Th(n)$ soltanto nel caso peggiore. È
proprio questa condivisione il risparmio che il keyword tree realizza rispetto a $z$
strutture separate.

L'altezza è la lunghezza del pattern più lungo. Il grado di un nodo è $\le \abs{\S}$ per la
proprietà (1) della definizione~\ref{def:keyword-tree}, ed è $\le z$ perché ogni sottoalbero figlio
contiene almeno una foglia e le foglie sono esattamente $z$.
\end{osservazione}

\begin{osservazione}[Nota implementativa]
Rappresentando i figli di un nodo con un array indicizzato su $\S$ si ha accesso $\Oh(1)$ al
figlio dato il carattere, ma spazio $\Oh(n\,\abs{\S})$; con liste, alberi bilanciati o
tabelle hash lo spazio resta $\Oh(n)$ e l'accesso costa $\Oh(\log\abs{\S})$ oppure $\Oh(1)$
atteso. Tutte le complessità «lineari» che seguono vanno intese nel modello ad accesso
$\Oh(1)$ ai figli.
\end{osservazione}

\section{I failure link}
\label{sec:failure-link}

Per utilizzare $K(\mathcal{P})$ durante la ricerca dobbiamo completarlo con i \emph{failure
link}, pensati per consentire al puntatore sul testo di avanzare \emph{sempre}: in caso di
mismatch non si ricomincia la discesa dalla radice riportando indietro il puntatore su $T$,
ma si salta al nodo la cui etichetta di cammino è il più lungo suffisso \emph{proprio} di
quanto già riconosciuto che sia ancora etichetta di un cammino dalla radice in
$K(\mathcal{P})$ --- equivalentemente, che sia ancora prefisso di almeno un pattern.

\begin{figure}[H]
\centering
% Idea del failure link: se la discesa si blocca in v si salta a n_v, il nodo la cui
% etichetta di cammino alfa e' il piu' lungo suffisso proprio di L(v) che sia ancora
% un cammino dalla radice.  I due tratti marcati alfa hanno la stessa lunghezza.
\begin{tikzpicture}[
    cammino/.style = {draw=graydark, line width=0.6pt, line cap=round, line join=round},
    marcato/.style = {draw=accent, line width=2.6pt, line cap=round, opacity=0.7},
    nodo/.style    = {circle, fill=graydark, inner sep=0pt, minimum size=3.6pt},
    flink/.style   = {draw=black, line width=0.7pt, dash pattern=on 3.5pt off 2.5pt,
                      -{Triangle[length=5.5pt,width=4.5pt]}},
    prosegue/.style= {draw=graydark, line width=0.6pt, dotted, line cap=round},
    et/.style      = {font=\footnotesize},
  ]

  % ---- sagoma di K(P) ---------------------------------------------------
  \fill[black!5, draw=gray!55, line width=0.4pt]
      (0,4.8) .. controls (2.02,4.2) and (3.70,2.49) .. (4.4,0.145)
              .. controls (2.2,-0.30) and (-2.2,-0.30) .. (-4.4,0.145)
              .. controls (-3.70,2.49) and (-2.02,4.2) .. cycle;

  % ---- nodi -------------------------------------------------------------
  \coordinate (r)  at ( 0.00, 4.80);
  \coordinate (u)  at (-1.02, 3.07);
  \coordinate (v)  at (-2.30, 0.90);
  \coordinate (nv) at ( 1.36, 2.70);

  % ---- cammini: i due tratti che portano alfa sono marcati ---------------
  \draw[marcato] (u) -- (v);
  \draw[marcato] (r) -- (nv);
  \draw[cammino] (r) -- (u) -- (v);
  \draw[cammino] (r) -- (nv);
  \draw[prosegue] (nv) -- (1.78,2.05);

  \node[nodo] at (r)  {};
  \node[nodo] at (u)  {};
  \node[nodo] at (v)  {};
  \node[nodo] at (nv) {};

  % ---- failure link -----------------------------------------------------
  \draw[flink] (v) -- (nv)
      node[et,midway,sloped,above=1pt] {\emph{failure link}};

  % ---- etichette --------------------------------------------------------
  \node[et,above=2pt]           at (r)  {$\mathrm{root}$};
  \node[et,anchor=east,xshift=-3pt]  at (u)  {$u$};
  \node[et,anchor=east,xshift=-3pt]  at (v)  {$v$};
  \node[et,anchor=west,xshift=4pt,yshift=3pt] at (nv) {$\nv{v}$};
  \node[et] at (-2.12, 2.24) {$\a$};
  \node[et] at ( 1.08, 4.00) {$\a$};
  \node[et] at ( 2.95, 0.70) {$K(\mathcal{P})$};
\end{tikzpicture}

\caption{Idea del failure link. Se la discesa si blocca nel nodo $v$, si salta al nodo
$\nv{v}$, la cui etichetta di cammino $\a$ è il più lungo suffisso proprio dell'etichetta di
cammino di $v$ che sia ancora un cammino dalla radice in $K(\mathcal{P})$: la porzione di
testo già riconosciuta resta allineata e il puntatore su $T$ non arretra.}
\label{fig:failure-link-idea}
\end{figure}

\subsection{Formalizzazione}

\scan{14}

\begin{definizione}[Etichetta di cammino, failure node]\label{def:ac-failure-node}
Dato un nodo $v$ di $K(\mathcal{P})$ definiamo:
\begin{itemize}
  \item $\L(v)$: la concatenazione delle etichette degli archi nel cammino
        $\mathrm{root}(K(\mathcal{P})) \rightsquigarrow v$;
  \item $\nv{v}$: il nodo di $K(\mathcal{P})$ tale che $\L(\nv{v})$ è il più lungo suffisso
        proprio di $\L(v)$ fra quelli che etichettano un cammino uscente dalla radice, cioè
        che sono della forma $\L(u)$ per qualche nodo $u$ di $K(\mathcal{P})$
        (equivalentemente: che sono prefissi di almeno un pattern di $\mathcal{P}$);
  \item $\mathit{lp}(v) \defeq \abs{\L(\nv{v})}$.
\end{itemize}
\end{definizione}

\begin{osservazione}
Attenzione al pedice, perché la lettera $n$ è ormai sovraccarica: in $n_i$ il pedice è
l'indice di un pattern e $n_i$ è una \emph{lunghezza}, in $\nv{v}$ il pedice è un nodo e
$\nv{v}$ è a sua volta un \emph{nodo} di $K(\mathcal{P})$.

Ogni arco di $K(\mathcal{P})$ porta esattamente un carattere, quindi
$\abs{\L(u)} = \mathrm{depth}(u)$ per ogni nodo $u$: dunque $\mathit{lp}(v)$ è la
\emph{profondità} di $\nv{v}$. Inoltre $\L(\nv{v})$ è un suffisso \emph{proprio} di $\L(v)$,
e perciò
\[
  \mathrm{depth}(\nv{v}) \;<\; \mathrm{depth}(v),
\]
fatto da cui dipenderà la terminazione di tutto ciò che segue. Per $z = 1$ il keyword tree
degenera in un cammino e $\mathit{lp}(v)$ è la lunghezza del più lungo bordo proprio del
prefisso già letto: la mappa $v \mapsto \nv{v}$ è la generalizzazione ad albero della
funzione di fallimento di KMP.
\end{osservazione}

\begin{lemma}[Buona definizione di $\nv{v}$]\label{lem:ac-nv}
Per ogni nodo $v \in K(\mathcal{P})$ con $v \neq \mathrm{root}(K(\mathcal{P}))$ esiste ed è
unico il nodo $\nv{v}$.
\end{lemma}

\begin{dimostrazione*}
Esercizio. \emph{N.B.} $\nv{v}$ può essere la radice.
\end{dimostrazione*}

\begin{osservazione}
L'enunciato va inteso per $v \neq \mathrm{root}(K(\mathcal{P}))$: infatti
$\L(\mathrm{root}) = \vuota$ non ha suffissi propri. Per la radice si adotta la convenzione
$\nv{\mathrm{root}} \defeq \mathrm{root}$ e $\mathit{lp}(\mathrm{root}) \defeq 0$.
\end{osservazione}

\begin{definizione}[Failure link]\label{def:ac-failure-link}
La coppia $\tuplev{v, \nv{v}}$ si chiama \emph{failure link}.
\end{definizione}

\subsection{Calcolo dei failure link}
\label{sec:calcolo-failure-link}

I failure link vanno calcolati contestualmente alla costruzione di $K(\mathcal{P})$, cioè
nella stessa fase di preprocessing, prima di toccare $T$: serve dunque un \emph{algoritmo
per determinare $\nv{v}$ dato $v$}. L'idea è ricorsiva e si legge sulla
figura~\ref{fig:ac-failure-link-passo}: se $v'$ è il padre di $v$ e $x$ è l'etichetta dell'arco
$v' \to v$ (cosicché $\L(v) = \L(v')\,x$), allora, noto $\nv{v'}$, basta guardare se
$\nv{v'}$ ha un arco uscente etichettato $x$.

\begin{figure}[H]
\centering
% Passo ricorsivo del calcolo dei failure link in Aho--Corasick.
% A sinistra il cammino root ~> u ~> v' -> v; a destra il cammino root ~> n_{v'} -> n_v.
% I due tratti marcati portano la stessa stringa alfa = L(n_{v'}); i due archi x sono
% paralleli e di uguale lunghezza, cosi' la regola si legge a colpo d'occhio.
\begin{tikzpicture}[
    cammino/.style = {draw=graydark, line width=0.6pt, line cap=round, line join=round},
    marcato/.style = {draw=accent, line width=2.6pt, line cap=round, opacity=0.7},
    tacca/.style   = {draw=black, line width=1.1pt, line cap=round},
    prosegue/.style= {draw=graydark, line width=0.6pt, dotted, line cap=round},
    noto/.style    = {draw=black, line width=0.7pt, dash pattern=on 3.5pt off 2.5pt,
                      -{Triangle[length=5.5pt,width=4.5pt]}},
    cercato/.style = {draw=gray, line width=0.7pt, dash pattern=on 1.6pt off 2pt,
                      -{Triangle[length=5pt,width=4pt]}},
    et/.style      = {font=\footnotesize},
    pic/.style     = {font=\scriptsize},
  ]

  % ---- sagoma di K(P) ---------------------------------------------------
  \fill[black!5, draw=gray!55, line width=0.4pt]
      (0,5.6) .. controls (2.30,4.90) and (4.20,2.90) .. (5.0,0.15)
              .. controls (2.5,-0.35) and (-2.5,-0.35) .. (-5.0,0.15)
              .. controls (-4.20,2.90) and (-2.30,4.90) .. cycle;

  % ---- nodi -------------------------------------------------------------
  \coordinate (r)   at ( 0.000, 5.60);
  \coordinate (u)   at (-1.350, 3.55);
  \coordinate (vp)  at (-2.850, 1.55);
  \coordinate (v)   at (-3.470, 0.55);
  \coordinate (np)  at ( 1.325, 3.48);
  \coordinate (nv)  at ( 1.945, 2.48);

  % ---- tratti marcati (la stessa stringa alfa) --------------------------
  \draw[marcato] (u) -- (vp);
  \draw[marcato] (r) -- (np);

  % ---- cammino di sinistra ---------------------------------------------
  \draw[cammino] (r) -- (u) -- (vp) -- (v);
  \draw[prosegue] (v) -- (-3.80,0.02);

  % ---- cammino di destra ------------------------------------------------
  \draw[cammino] (r) -- (np) -- (nv);
  \draw[cammino] (np) -- (2.16,3.14);          % moncone dell'arco y
  \draw[prosegue] (nv) -- (2.26,1.98);

  % ---- tacche trasversali sui nodi --------------------------------------
  \draw[tacca] ($(u)!3pt!90:(vp)$)   -- ($(u)!3pt!-90:(vp)$);
  \draw[tacca] ($(vp)!4pt!90:(v)$)   -- ($(vp)!4pt!-90:(v)$);
  \draw[tacca] ($(v)!4pt!90:(vp)$)   -- ($(v)!4pt!-90:(vp)$);
  \draw[tacca] ($(np)!4pt!90:(nv)$)  -- ($(np)!4pt!-90:(nv)$);
  \draw[tacca] ($(nv)!4pt!90:(np)$)  -- ($(nv)!4pt!-90:(np)$);
  \fill[graydark] (r) circle (1.8pt);

  % ---- failure link -----------------------------------------------------
  \draw[noto]    (vp) -- (np)
      node[pic,pos=0.56,sloped,above=1pt] {\emph{failure link} (noto)};
  \draw[cercato] (v) -- (nv)
      node[pic,pos=0.50,sloped,below=1pt,text=gray] {\emph{failure link} cercato};

  % ---- etichette del cammino di sinistra --------------------------------
  \node[et,above=3pt] at (r) {$\mathrm{root}$};
  \node[et,anchor=east,xshift=-4pt] at (u)  {$u$};
  \node[et,anchor=east,xshift=-5pt,yshift=4pt] at (vp) {$v'$};
  \node[et,anchor=east,xshift=-5pt,yshift=-3pt] at (v) {$v$};
  \node[et,anchor=east] at (-3.44,1.14) {$x$};
  \node[et] at (-2.52,2.74) {$\a$};

  % ---- etichette del cammino di destra ----------------------------------
  \node[et,anchor=west] at (1.60,3.70) {$\nv{v'}$};
  \node[et,anchor=west,xshift=3pt,yshift=-2pt] at (nv) {$\nv{v}$};
  \node[et,anchor=east] at (1.52,2.84) {$x$};
  \node[et,anchor=west] at (2.24,3.06) {$y \neq x$};
  \node[et] at (1.06,4.76) {$\a$};

  \node[et] at (3.55,0.60) {$K(\mathcal{P})$};
\end{tikzpicture}

\caption{Passo ricorsivo del calcolo dei failure link. Il tratto marcato sul cammino di
sinistra (da $u$ a $v'$) e il cammino marcato di destra (da $\mathrm{root}$ a $\nv{v'}$)
portano la stessa stringa $\a$: è la relazione $\L(\nv{v'}) = \a$, il più lungo suffisso
proprio di $\L(v')$ che sia anche prefisso di un pattern. Se $\nv{v'}$ possiede un arco
uscente etichettato $x$, il nodo così raggiunto è $\nv{v}$.}
\label{fig:ac-failure-link-passo}
\end{figure}

\scan{15}
Svolgiamo la ricorsione, procedendo sulla profondità dei nodi.

\paragraph{Caso base.}
$v = \mathrm{root}$: banale, si pone $\nv{\mathrm{root}} = \mathrm{root}$. Banale è anche il
caso dei figli della radice: se $\mathrm{depth}(v) = 1$ allora $\L(v)$ è un singolo
carattere e il suo unico suffisso proprio è $\vuota = \L(\mathrm{root})$, dunque
$\nv{v} = \mathrm{root}$.

\paragraph{Passo induttivo.}
Sia ora $\mathrm{depth}(v) \ge 2$: i nodi di profondità $\le 1$ sono già coperti dal caso
base, e vanno esclusi qui perché per essi il ragionamento che segue darebbe la risposta
sbagliata (se $v' = \mathrm{root}$ allora $\nv{v'} = \mathrm{root}$, che un figlio raggiunto
dall'arco etichettato $x$ ce l'ha --- è $v$ stesso --- e si concluderebbe $\nv{v} = v$
anziché $\nv{v} = \mathrm{root}$).
Se devo determinare $\nv{v}$, sia $v' = \mathrm{parent}(v)$ e sia $x$ l'etichetta dell'arco
$\tuplev{v', v}$, cosicché $\L(v) = \L(v')\,x$. Poiché
\[
  \mathrm{depth}(v') < \mathrm{depth}(v)
  \quad\Longrightarrow\quad
  \text{ho già } \nv{v'},
\]
posso ragionare per casi:
\begin{itemize}
  \item se $\nv{v'}$ ha un figlio raggiunto da un arco etichettato $x$, allora $\nv{v}$ è
        quel figlio;
  \item altrimenti passo a $\nv{\nv{v'}}$ e determino se \emph{questo} ha un figlio
        etichettato $x$, e così via\dots
\end{itemize}
Si risale cioè la catena dei failure link
\[
  \nv{v'},\quad \nv{\nv{v'}},\quad \nv{\nv{\nv{v'}}},\quad \dots
\]
finché o si trova un nodo con un arco uscente etichettato $x$ --- e $\nv{v}$ è il figlio così
raggiunto --- oppure si arriva alla radice senza trovarlo, e in tal caso
$\nv{v} = \mathrm{root}$. La risalita termina sempre, perché lungo la catena la profondità
decresce strettamente: l'algoritmo non può andare in loop.

\paragraph{Ordine di elaborazione.}
Perché la ricorsione sia ben fondata i nodi vanno elaborati per profondità crescente, cioè
con una visita per livelli: quando si calcola $\nv{v}$ il valore $\nv{v'}$ del padre deve
essere già disponibile.

\begin{algorithm}[H]
\caption{Calcolo dei failure link di $K(\mathcal{P})$}
\label{alg:failure-link}
\begin{algorithmic}[1]
\Require il keyword tree $K(\mathcal{P})$
\Ensure $\nv{v}$ e $\mathit{lp}(v)$ per ogni nodo $v$ di $K(\mathcal{P})$
\State $\nv{\mathrm{root}} \gets \mathrm{root}$
\State $\mathit{lp}(\mathrm{root}) \gets 0$
\For{ogni figlio $u$ della radice}
  \State $\nv{u} \gets \mathrm{root}$
  \State $\mathit{lp}(u) \gets 0$
\EndFor
\For{$d \gets 2, \dots, \text{altezza di } K(\mathcal{P})$}
  \For{ogni nodo $v$ con $\mathrm{depth}(v) = d$}
    \State $v' \gets \mathrm{parent}(v)$
    \State $x \gets$ etichetta dell'arco $\tuplev{v', v}$
    \State $u \gets \nv{v'}$
    \While{$u \neq \mathrm{root}$ e $u$ non ha archi uscenti etichettati $x$}
      \State $u \gets \nv{u}$
    \EndWhile
    \If{$u$ ha un arco uscente $\tuplev{u, u'}$ etichettato $x$}
      \State $\nv{v} \gets u'$
    \Else
      \State $\nv{v} \gets \mathrm{root}$
    \EndIf
    \State $\mathit{lp}(v) \gets \mathrm{depth}(\nv{v})$
  \EndFor
\EndFor
\end{algorithmic}
\end{algorithm}

\paragraph{Complessità.}
Il numero di nodi analizzati per determinare un singolo $\nv{v}$ è
$\le \mathrm{depth}(v)$: da qui, mediante un'\textbf{analisi ammortizzata}, si ottiene una
complessità \textbf{lineare}.

\begin{osservazione}
Il conto per singolo nodo darebbe soltanto il limite pessimistico
$\Oh\bigl(n \cdot \max_i n_i\bigr)$. L'analisi ammortizzata va invece fatta lungo un intero
cammino radice--foglia, cioè lungo un pattern $P_i$: usando come potenziale la profondità
del candidato, $\mathit{lp}(\cdot) \ge 0$, si osserva che passando da un nodo al figlio
successivo essa cresce al più di $1$ (si aggiunge un solo carattere), mentre ogni iterazione
della risalita lungo la catena dei failure link la fa calare di almeno $1$. Il numero
complessivo di risalite lungo il cammino è dunque $\le n_i$: al peggio si paga il doppio dei
passi. Sommando su tutti i pattern,
\[
  T_{\text{preprocessing}} \;=\; \Oh\Bigl(\sum_{i=1}^{z} n_i\Bigr) \;=\; \Oh(n),
\]
nel modello in cui l'accesso al figlio dato il carattere costa $\Oh(1)$.
\end{osservazione}

% RICOSTRUZIONE (in fase di cucitura, non presente nel quaderno): il manoscritto non
% formalizza mai la fase di ricerca.  Le posizioni l e c, l'invariante L(v) = T[l..c-1],
% l'aggiornamento l <- c - lp(v), v <- n_v e la complessita' O(|T|+k) sono ricostruiti da
% Gusfield sec. 3.4 e da L04_sintesi sec. 12.  Senza questa sezione il capitolo definirebbe
% i failure link senza mai usarli.  Da registrare in NOTE-EDITORIALI.md (Ricostruzioni).
\section{La fase di ricerca}
\label{sec:ricerca-ac}

Il preprocessing ha prodotto $K(\mathcal{P})$ con tutti i suoi failure link; la scansione di
$T$ procede ora mantenendo un nodo corrente $v$ di $K(\mathcal{P})$, una posizione corrente
$c$ in $T$ e la posizione $\ell$ in cui inizia l'allineamento corrente, con l'invariante
$\L(v) = \sub{T}{\ell}{c-1}$. Finché da $v$ esce un arco etichettato $T[c]$ si scende
lungo quell'arco e si incrementa $c$; se si raggiunge una foglia si riporta l'occorrenza del
pattern corrispondente. In caso di mismatch si segue il failure link, ponendo nell'ordine
\[
  \ell \;\gets\; c - \mathit{lp}(v), \qquad v \;\gets\; \nv{v},
\]
senza mai decrementare $c$ (se il blocco avviene già sulla radice si incrementa
semplicemente $c$): è esattamente per questo che i failure link sono stati introdotti.
Poiché $c$ non arretra mai e a ogni salto $\ell$ cresce strettamente, la scansione costa
$\Oh(\abs{T} + k)$.

\section{Scarico dell'ipotesi: gli output link}
\label{sec:output-link}

\scan{15}

\begin{domanda}
Che succede se una stringa è \emph{fattore} di un'altra stringa?
\end{domanda}

\begin{esempio}
Si consideri
\[
  \mathcal{P} = \set{\str{acatt},\ \str{ca}},
  \qquad
  T = \str{acatg}.
\]
Qui $\str{ca}$ è sottostringa di $\str{acatt}$ (posizioni $2$--$3$): l'ipotesi semplificativa
«nessun $P_i$ è sottostringa di un altro $P_j$» è violata.
\end{esempio}

\begin{figure}[H]
\centering
% Keyword tree di P = {acatt, ca}.  La scansione di T = acatg scende lungo il ramo
% sinistro (a, c, a, t) e si blocca sulla g: la foglia di ca non viene mai raggiunta
% e l'occorrenza di ca in T[2..3] va perduta.
\begin{tikzpicture}[
    arco/.style   = {draw=graydark, line width=0.55pt, line cap=round},
    nodo/.style   = {circle, fill=black, inner sep=0pt, minimum size=3.2pt},
    car/.style    = {font=\footnotesize\ttfamily, text=graydark},
    et/.style     = {font=\footnotesize},
    lieve/.style  = {font=\scriptsize\ttfamily, text=gray},
    scan/.style   = {draw=accent, line width=2.8pt, line cap=round, line join=round,
                     opacity=0.55},
    chev/.style   = {draw=accent!65!black, line width=0.7pt, line cap=round,
                     -{Stealth[length=4.5pt,width=3.4pt]}},
  ]

  % ---- nodi -------------------------------------------------------------
  \coordinate (r)  at ( 0.0,  0.0);
  \coordinate (u1) at ( 1.2, -1.1);
  \coordinate (u2) at ( 2.2, -2.2);
  \coordinate (v1) at (-1.0, -1.1);
  \coordinate (v2) at (-1.9, -2.2);
  \coordinate (v3) at (-2.8, -3.3);
  \coordinate (v4) at (-3.7, -4.4);
  \coordinate (v5) at (-4.6, -5.5);

  % ---- cammino di scansione (evidenziato) -------------------------------
  \draw[scan, rounded corners=8pt]
      (0.15,-0.35) -- (-0.55,-1.05) -- (-1.45,-2.15) -- (-2.35,-3.25) -- (-3.45,-4.55);
  \draw[accent!65!black, line width=0.9pt, line cap=round,
        -{Triangle[length=5pt,width=4.5pt]}] (-3.50,-4.55) -- (-3.50,-5.58);
  \foreach \p in {{-0.55,-1.05},{-1.45,-2.15},{-2.35,-3.25}}
      \draw[chev] ($(\p)+(51:0.22)$) -- ($(\p)+(231:0.16)$);

  % ---- archi ------------------------------------------------------------
  \draw[arco] (r) -- (v1) -- (v2) -- (v3) -- (v4) -- (v5);
  \draw[arco] (r) -- (u1) -- (u2);

  \foreach \n in {r,u1,u2,v1,v2,v3,v4,v5} \node[nodo] at (\n) {};

  % ---- etichette degli archi -------------------------------------------
  \node[car] at (-0.74,-0.34) {a};
  \node[car] at (-1.70,-1.45) {c};
  \node[car] at (-2.60,-2.55) {a};
  \node[car] at (-3.50,-3.65) {t};
  \node[car] at (-4.40,-4.75) {t};
  \node[car] at ( 0.82,-0.32) {c};
  \node[car] at ( 1.94,-1.44) {a};

  % ---- fine dei due pattern --------------------------------------------
  \node[lieve,anchor=north] at ($(v5)+(0,-0.12)$) {acatt};
  \node[lieve,anchor=west]  at ($(u2)+(0.16,-0.06)$) {ca};

  % ---- punto di fallimento ---------------------------------------------
  \draw[accent!65!black, line width=1.5pt, line cap=round]
      (-3.74,-4.88) -- (-3.26,-5.36)  (-3.74,-5.36) -- (-3.26,-4.88);

  \node[et, anchor=north west, align=left] at (-3.08,-4.94)
      {ho perso l'occorrenza\\ di \str{ca}!};
\end{tikzpicture}

\caption{Keyword tree di $\mathcal{P} = \set{\str{acatt}, \str{ca}}$. Scandendo
$T = \str{acatg}$ la discesa segue gli archi $\str{a}, \str{c}, \str{a}, \str{t}$ e si
arresta sul carattere $\str{g}$ ($\times$): la foglia di \str{ca} non viene mai visitata e
l'occorrenza di \str{ca} in $\sub{T}{2}{3}$ va perduta.}
\label{fig:ac-output-link}
\end{figure}

Percorrendo il cammino di \str{acatt} si passa «sopra» un'occorrenza di \str{ca} senza mai
visitare la foglia corrispondente: ho perso l'occorrenza di \str{ca}!

\paragraph{Soluzione.}
Aggiungo a ogni nodo un \emph{output link} (equivalentemente, una \emph{output list}) di
occorrenze da riportare, quelle dovute alle sottostringhe. Per ogni nodo $v$ si pone
\[
  \out(v) \;=\;
  \begin{cases}
    \nv{v} & \text{se } \L(\nv{v}) \in \mathcal{P},\\[2pt]
    \out(\nv{v}) & \text{altrimenti,}
  \end{cases}
\]
con $\out(v)$ indefinito se la catena raggiunge la radice senza incontrare la fine di un
pattern. Arrivati in $v$ durante la ricerca si percorre la catena degli output link
riportando tutte le occorrenze che terminano nella posizione corrente. Poiché $\out(v)$ si
calcola in $\Oh(1)$ per nodo \emph{durante} la stessa visita per livelli che calcola i
failure link (algoritmo~\ref{alg:failure-link}), il preprocessing resta $\Oh(n)$ e la ricerca
resta $\Oh(\abs{T} + k)$. Scaricata l'ipotesi, va anche marcato come «fine di un pattern»
ogni nodo, interno o foglia, la cui etichetta di cammino appartiene a $\mathcal{P}$.

\section{Bilancio}

Mettendo insieme le due fasi, l'algoritmo di Aho--Corasick costa
\[
  T_{\text{preprocessing}} = \Oh(n),
  \qquad
  T_{\text{ricerca}} = \Oh(\abs{T} + k),
  \qquad
  T_{\text{totale}} = \Oh\bigl(n + \abs{T} + k\bigr),
\]
da confrontare con l'$\Oh(z\,\abs{T} + n)$ della soluzione naive: il fattore $z$ è sparito e
il testo viene scandito una volta sola, che era esattamente la sfida posta nella
\autoref{sec:exact-set-matching}.
